\documentclass{beamer}
\usetheme{Madrid}

% --- PREAMBLE ---
\usepackage[utf8]{inputenc}
\usepackage[T1]{fontenc}
\usepackage{amsmath}

% Paket untuk Teorema
\usepackage{tcolorbox}
\tcbuselibrary{skins}
\usepackage{ifthen}
\usepackage{xparse} % WAJIB untuk argumen opsional ganda

% Paket untuk referensi cerdas
\usepackage{cleveref}

% --- DEFINISI COUNTER, CLEVEREF, DAN ENVIRONMENT ---
\newcounter{nomorTeorema}
\renewcommand{\thenomorTeorema}{\arabic{nomorTeorema}}
\crefname{nomorTeorema}{Teorema}{Teorema-teorema} % Pengaturan untuk \cref

% Definisi Environment Final (Antipeluru)
% Argumen #1: Judul Tambahan (Opsional)
% Argumen #2: Nama Label (Opsional)
\NewDocumentEnvironment{Teorema}{ O{} O{} }
  { % \begin
    \only<1>{
      \stepcounter{nomorTeorema}
      \ifstrempty{#2}{}{
        \label{#2}
      }
    }
    \begin{tcolorbox}[
      enhanced, colback=teal!10, colframe=teal!80!black,
      fonttitle=\bfseries, coltitle=white, colbacktitle=teal!70!black,
      title={Teorema \thenomorTeorema\ifstrempty{#1}{}{:\ #1}}
    ]
    \setbeamertemplate{enumerate items}[default]
    \setlength{\leftmargini}{6mm}
  }
  { % \end
    \end{tcolorbox}
  }
% -------------------------------------------

\begin{document}

\begin{frame}{Uji Coba Penomoran dan Referensi Final}

  % Teorema dengan judul dan label
  \begin{Teorema}[Pythagoras][thm:pythagoras]
    Untuk segitiga siku-siku, $a^2 + b^2 = c^2$.
  \end{Teorema}

  \pause % Jeda untuk memastikan overlay tidak merusak nomor

  % Teorema tanpa judul, hanya label
  \begin{Teorema}[][thm:lain]
    Ini adalah teorema kedua. Referensinya harus menunjuk ke nomor 2.
  \end{Teorema}

  \vfill

  % Merujuk teorema-teorema tersebut
  Nomor referensi untuk \cref{thm:pythagoras} sekarang dijamin benar.
  Begitu pula dengan referensi untuk \cref{thm:lain}.

  \begin{alertblock}{Penting!}
    Jika nomor masih tampak salah, \textbf{kompilasi ulang dokumen Anda sekali lagi}. LaTeX terkadang membutuhkan dua kali kompilasi untuk memperbarui semua referensi dengan benar.
  \end{alertblock}

\end{frame}

\end{document}